\section{Introduction}
\label{sec:introduction}

Confluence is rewriting theory's cleanest order-independence principle: if a state reduces in two different ways, the divergence should be joinable \cite{BaaderNipkow1998,Huet1980,DershowitzJouannaud1990,Terese2003}. Discrete transport geometry packages the same intuition in a different language: transport should be independent of route, and the obstruction is holonomy \cite{Meneses2021,BaezMuniain1994}. This paper argues that their local cores coincide once rewriting is given its missing two-dimensional cells. The relevant object is not the bare reduction graph. It is a reduction 2-complex whose $0$-cells are states, whose $1$-cells are one-step rewrites, and whose generating $2$-cells are local peaks or, in structured rewriting settings, reachable critical pairs \cite{Mimram2023,GuiraudMalbos2018,GuiraudMalbos2012}.

The first temptation is also the first mistake. If one tries to read flatness directly from the bare reduction graph, one obtains surrogates that are vacuous or misleading in the terminating settings studied here. \Cref{sec:bare-graph-not-enough} shows this explicitly: directed-loop flatness collapses to acyclicity on the curated terminating examples, while ordinary undirected graph cycle structure is not diagnostic for confluence. The negative result is therefore constructive. It tells us where the geometric data must live: not on graph loops alone, but on elementary branching cells.

Once those missing $2$-cells are restored, the local theorem-level core becomes clean. \Cref{sec:core-theorem} proves that elementary flatness is equivalent to local confluence. The paper then shows computationally that the corrected local/global picture survives exhaustive finite-ARS stress tests in the terminating regime, curated bounded string-rewriting examples, and a curated left-linear term-rewriting bridge with genuine nonroot overlap. In those structured settings, the relevant local generators are reachable critical pairs, and the before/after completion examples show that defect elimination can be read as holonomy elimination.

This paper is also a mathematics-facing contribution to the Six Birds line of work. The general closure vocabulary comes from \cite{Tsiokos2026SixBirdsFoundations}, while the treatment of route mismatch as a disciplined mathematical defect was developed earlier in \cite{Tsiokos2026CountStone}. The present paper does not repeat that general background. Instead, it uses that vocabulary to isolate a specific claim about rewriting: the local obstructions to confluence are carried by the generating $2$-cells on which elementary holonomy is tested.

The paper is organized around three levels of claim. First, there is a proved core: elementary flatness is equivalent to local confluence, and confluence plus normalization yields a unique normal form from a chosen start state. Second, there is a strong computational and structural bridge: exhaustive finite-ARS audits, curated string and term-rewriting examples, and explicit reduction 2-complex artifacts all support the corrected local geometric picture. Third, there is a deferred theorem package: the full finite terminating left-linear term-rewriting formulation in the exact target form remains to be written as a general theorem rather than a curated bridge. \Cref{sec:proved-supported-deferred} makes this boundary explicit.

\paragraph{Contributions.}
This paper makes four contributions.
\begin{enumerate}
    \item It identifies the correct geometric object for rewriting: a reduction 2-complex rather than a bare reduction graph.
    \item It proves the local formal core that elementary flatness is exactly local confluence (\Cref{thm:elementary-flatness-iff-local-confluence}), together with the supporting uniqueness-of-normal-form fact under confluence and normalization (\Cref{prop:confluence-normalization-unique-normal-form}).
    \item It shows computationally that the corrected local/global picture survives exhaustive finite-ARS stress tests and curated string and left-linear term-rewriting examples, while graph-only flatness surrogates fail (\Cref{sec:bare-graph-not-enough,sec:computational-bridge}).
    \item It gives an explicit operational reading of completion as holonomy elimination in before/after string and term-rewriting pairs (\Cref{sec:completion-holonomy}).
\end{enumerate}

The rest of the paper proceeds as follows. \Cref{sec:preliminaries-scope} fixes the paper's scope and its minimal Six Birds placement. \Cref{sec:bare-graph-not-enough} rules out the wrong graph-only formulation. \Cref{sec:reduction-two-complex} introduces the reduction 2-complex and elementary holonomy language, and \Cref{sec:core-theorem} proves the theorem-level core. \Cref{sec:critical-pairs-curvature} identifies reachable critical pairs as the structured rewriting generators of local curvature, \Cref{sec:computational-bridge} records the finite-ARS, string, and term-rewriting evidence base, and \Cref{sec:completion-holonomy} packages the before/after examples as holonomy-elimination case studies. \Cref{sec:proved-supported-deferred} then states explicitly what is proved, what is strongly supported, and what remains final theorem-writing work.
